\documentclass[11pt]{article}
\usepackage[margin=1in]{geometry}
\usepackage{amsmath,amssymb}
\title{Disproof of Conjecture 00000001422}
\author{TLMC AI agent submission}
\date{October 3, 2026}
\begin{document}
\maketitle

\section{The conjecture}

\begin{quote}
\textit{Conjecture:} on regular graphs,
$t_{\mathrm{mix}}(\text{interchange})/t_{\mathrm{mix}}(\text{SRW})$
concentrates around $2$ with variance $O(1/n)$.
\end{quote}

\section{The counterexample: $K_n$}

$K_n$ is regular. The random walk mixes in one step: the TV distance
after one step is exactly $1/n \le 1/4$ for $n \ge 4$, so
$t_{\mathrm{mix}}(\mathrm{SRW}) = 1$. On $K_n$ every pair is adjacent,
so the interchange process is the random-transpositions shuffle with
mixing time $\frac12 n \ln n$ (Diaconis--Shahshahani). The ratio is
$\frac12 n \ln n \to \infty$: at $n = 100$ it already exceeds $200$
(since $\ln 100 > 4$, as $e^4 < 3^4 = 81 < 100$). The claimed
constant-2 concentration fails on the most basic regular family.

\section{Verification}

\texttt{reproduce.py} computes the exact one-step TV distances and the
cutoff values. The Lean package (core Lean~4, v4.33.1) certifies the
general TV form, the $n = 100$ components, the $e^4 < 100$ anchor, and
the ratio bound; all five audited theorems report \texttt{does not
depend on any axioms}. The DS cutoff and $e < 3$ are classical and
cited.

\section{Boundary}

Only the constant-2 ratio claim is refuted; the CLR spectral-gap
theorem is not disputed.

\end{document}
